% !TEX program = xelatex
% Bản dùng HUST Beamer Template (SoICT) -- theme=blue, 16:9.
% Biên dịch bằng XeLaTeX (hỗ trợ tiếng Việt):
%   tectonic -X compile main.tex      hoặc      latexmk -xelatex main.tex
% Trên Overleaf: Menu -> Settings -> Compiler -> XeLaTeX.
\documentclass[aspectratio=169,10pt]{beamer}
\usetheme[theme=blue,logo=hustwithtext,assetext=pdf]{HUST}
\setbeamertemplate{navigation symbols}{}

% Màu chữ/ký hiệu theo nhận diện HUST
\setbeamercolor{structure}{fg=HUSTBlue}
\setbeamercolor{alerted text}{fg=HUSTRed}
\setbeamercolor{block title}{bg=HUSTBlue,fg=white}
\setbeamercolor{block body}{bg=HUSTBlue!8}
\setbeamercolor{block title alerted}{bg=HUSTRed,fg=white}
\setbeamercolor{block body alerted}{bg=HUSTRed!8}

% Số trang ở góc dưới phải, ngay trên thanh đỏ của footer
\setbeamertemplate{footline}{%
  \hfill{\scriptsize\color{HUSTBlue}\insertframenumber\,/\,\inserttotalframenumber}%
  \hspace{0.02\paperwidth}\vspace{0.075\paperheight}}

% ---------- Ngôn ngữ & font ----------
\usepackage{fontspec}          % font mặc định Latin Modern, đủ dấu tiếng Việt
\usepackage{polyglossia}
\setmainlanguage{vietnamese}
\setotherlanguage{english}

\usepackage{amsmath,amssymb}
\usepackage{booktabs,array}
\usepackage{pifont}
\usetikzlibrary{positioning,arrows.meta,fit}
\colorlet{mainblue}{HUSTBlue}

\newcommand{\yes}{\ding{51}}
\newcommand{\no}{\ding{55}}
\tikzset{
  stage/.style={draw=HUSTBlue, fill=HUSTBlue!8, rounded corners=3pt,
    align=center, font=\small, minimum height=1.1cm, text width=2.3cm},
  arrow/.style={-{Stealth[length=5pt]}, thick, HUSTBlue}
}

% Đặt nội dung tại vị trí cố định trên trang (dùng cho slide bìa / cảm ơn)
\newcommand{\placecontent}[4]{%
  \tikz[remember picture,overlay]
    \node[anchor=north west]
      at ([xshift=#1,yshift=-#2]current page.north west)
      {\parbox{#3}{#4}};
}

% Mỗi phần mở đầu bằng slide chuyển phần trên nền theme của HUST
\AtBeginSection[]{%
  {\HUSTUseBackground{theme}
  \begin{frame}[plain]
    \centering
    {\color{white}\scalebox{2}{\Huge\bfseries 0\thesection}\par}
    \vspace{10pt}
    {\color{white}\Large\bfseries\insertsectionhead\par}
  \end{frame}}}

\title{Multi-Agent Systems \& Collaborative Reasoning}
\subtitle{Khảo sát, phân tích và đánh giá 8 nghiên cứu năm 2026}
\author{Trần Xuân Hương · Vũ Tùng Dương · Phạm Việt Dũng · Lê Văn Đức}

\begin{document}

\HUSTInsertBrandSlide

{\HUSTUseBackground{onelove}
\begin{frame}[plain]
  \HUSTCornerImage[0.12]{assets/logo/03.pdf}
  \placecontent{0.55cm}{0.20\paperheight}{0.62\paperwidth}{
    \raggedright\color{HUSTRed}\bfseries\fontsize{22pt}{28pt}\selectfont
    \inserttitle
  }
  \placecontent{0.55cm}{0.43\paperheight}{0.60\paperwidth}{
    \raggedright\color{HUSTBlue}
    {\fontsize{12pt}{15pt}\selectfont \insertsubtitle\par}
    \vspace{14pt}
    \small
    \begin{minipage}[t]{0.27\paperwidth}
      \textbf{Nhóm thực hiện}\\[3pt]
      Trần Xuân Hương\\
      Vũ Tùng Dương\\
      Phạm Việt Dũng\\
      Lê Văn Đức
    \end{minipage}%
    \hspace{0.03\paperwidth}%
    \begin{minipage}[t]{0.28\paperwidth}
      \textbf{Giảng viên hướng dẫn}\\[3pt]
      PGS.TS. Lê Thanh Hương\\[8pt]
      \textbf{Học phần}\\[3pt]
      IT4772 -- Xử lý ngôn ngữ tự nhiên\\
      Báo cáo Bài tập lớn
    \end{minipage}
  }
\end{frame}
}

\begin{frame}{Nội dung trình bày}
  \tableofcontents
\end{frame}

% ============================================================
\section{Giới thiệu}
% ============================================================

\begin{frame}{LLM-based MAS: nhiều agent cùng suy luận}
  \begin{columns}[T]
    \begin{column}{0.5\textwidth}
      \textbf{Ý tưởng}
      \begin{itemize}
        \item Mỗi agent = backbone LLM + role prompt + (tùy chọn) tool.
        \item Các agent trao đổi message để cùng giải bài toán: planner, coder, critic, verifier…
        \item Framework tiêu biểu: CAMEL, AutoGen, MetaGPT, multi-agent debate.
      \end{itemize}
      \vspace{4pt}
      \textbf{Câu hỏi nghiên cứu năm 2026}\\
      Không còn là “cộng tác có giúp không?” mà là:
      hiệu quả chi phí · luồng thông tin · an toàn · ứng dụng chuyên ngành.
    \end{column}
    \begin{column}{0.46\textwidth}
      \begin{block}{MAS = communication graph}
        Node là agent; cạnh $v_j \to v_i$: output của $v_j$ làm context cho $v_i$.
        \[ r_i = \mathrm{LLM}(p_i,\, q,\, O_i) \]
        Topology có thể tối ưu (GPTSwarm, G-Designer) hoặc cắt tỉa (AgentPrune).
      \end{block}
      \begin{alertblock}{Vì sao topology quan trọng?}
        Vừa là đòn bẩy hiệu năng, vừa là kênh lan truyền lỗi, vừa là bề mặt tấn công.
      \end{alertblock}
    \end{column}
  \end{columns}
\end{frame}

\begin{frame}{Bốn thách thức chính của MAS}
  \begin{columns}[T]
    \begin{column}{0.48\textwidth}
      \begin{block}{C1 -- Chi phí}
        Nhiều lần gọi LLM, context dài $\to$ token, năng lượng GPU, latency tăng nhanh.\\
        $\Rightarrow$ ProgRouter, DAIN
      \end{block}
      \begin{block}{C2 -- Quá tải context}
        “Lost in the middle”; compaction hay pruning đều làm mất thông tin.\\
        $\Rightarrow$ AGENT-RADAR
      \end{block}
    \end{column}
    \begin{column}{0.48\textwidth}
      \begin{block}{C3 -- Lan truyền lỗi}
        Hallucination hoặc misinformation lan theo cạnh và bị khuếch đại.\\
        $\Rightarrow$ HalluProp, TIDE
      \end{block}
      \begin{block}{C4 -- Bảo mật topology}
        Topology là tài sản trí tuệ; lộ topology $\to$ tấn công có chủ đích.\\
        $\Rightarrow$ MIRAGE
      \end{block}
    \end{column}
  \end{columns}
\end{frame}

\begin{frame}{Tám paper được khảo sát}
  \begin{columns}[T]
    \begin{column}{0.68\textwidth}
      \scriptsize
      \begin{tabular}{@{}c l >{\raggedright\arraybackslash}p{2.9cm} c l@{}}
        \toprule
        \# & Paper & Chủ đề chính & Nộp & Nhóm \\
        \midrule
        1 & MIRAGE      & Che giấu topology bằng phantom structure    & 29/09 & Safety \\
        2 & MAC-RRG     & Multi-agent cho radiology report generation & 20/09 & Application \\
        3 & ProgRouter  & Routing LLM theo tiến độ, có ngân sách      & 26/08 & Orchestration \\
        4 & TIDE        & Misinformation khi sự thật bị phân tán      & 04/08 & Safety \\
        5 & HalluProp   & Dự báo rủi ro hallucination pre-hoc         & 29/07 & Orchestration \\
        6 & UrbanAgent  & Urban region profiling (KDD’26)             & 15/07 & Application \\
        7 & DAIN        & Dynamic agent interaction network           & 29/06 & Communication \\
        8 & AGENT-RADAR & Attention steering cho MAS                  & 28/05 & Communication \\
        \bottomrule
      \end{tabular}
    \end{column}
    \begin{column}{0.28\textwidth}
      \textbf{Cách chọn bài}
      \begin{itemize}
        \item Nguồn: arXiv.
        \item Abstract chứa “multi-agent” và “collaborative reasoning”.
        \item 8 bài mới nhất (05--09/2026).
      \end{itemize}
      \vspace{4pt}
      \textbf{Lưu ý:} 7/8 là preprint; chỉ UrbanAgent đã được nhận ở KDD 2026.
    \end{column}
  \end{columns}
\end{frame}

\begin{frame}{Taxonomy: bốn hướng nghiên cứu}
  \begin{columns}[T]
    \begin{column}{0.48\textwidth}
      \begin{block}{Safety \& Robustness}
        Hệ thống có chống chịu được tấn công và thông tin sai?\\
        \textbf{\#1 MIRAGE · \#4 TIDE}
      \end{block}
      \begin{block}{Communication \& Interaction}
        Agent trao đổi gì, với ai, mức độ nào?\\
        \textbf{\#7 DAIN · \#8 AGENT-RADAR}
      \end{block}
    \end{column}
    \begin{column}{0.48\textwidth}
      \begin{block}{Orchestration \& Control}
        Ai chạy, chạy bằng model nào, agent nào rủi ro?\\
        \textbf{\#3 ProgRouter · \#5 HalluProp}
      \end{block}
      \begin{block}{Domain Applications}
        Áp dụng cộng tác vào bài toán chuyên ngành.\\
        \textbf{\#2 MAC-RRG · \#6 UrbanAgent}
      \end{block}
    \end{column}
  \end{columns}
  \vspace{4pt}
  {\small Phân loại theo mục tiêu chính; nhiều bài thuộc nhiều nhóm
  (VD: HalluProp vừa là orchestration vừa là safety).}
\end{frame}

\begin{frame}{Mỗi phương pháp can thiệp ở đâu trong lifecycle của MAS}
  \centering
  \begin{tikzpicture}[node distance=0.35cm, stage/.append style={text width=2.0cm}]
    \node[stage] (s1) {Thiết kế role\\\& graph};
    \node[stage, right=of s1] (s2) {Điều phối\\\& kích hoạt};
    \node[stage, right=of s2] (s3) {Giao tiếp\\\& suy luận};
    \node[stage, right=of s3] (s4) {Tổng hợp\\\& phơi bày};
    \node[stage, right=of s4] (s5) {Đánh giá\\\& kiểm toán};
    \foreach \a/\b in {s1/s2,s2/s3,s3/s4,s4/s5} \draw[arrow] (\a) -- (\b);
    \foreach \s/\list in {
        s1/{HalluProp (pre-hoc)},
        s2/{ProgRouter\\DAIN (gating)},
        s3/{AGENT-RADAR\\DAIN (graph)\\UrbanAgent\\MAC-RRG},
        s4/{MIRAGE},
        s5/{TIDE}}
      \node[below=0.3cm of \s, align=center, font=\small, anchor=north] {\list};
  \end{tikzpicture}
\end{frame}

% ============================================================
\section{Safety \& Robustness}
% ============================================================

\begin{frame}{MIRAGE: che giấu topology bằng phantom structure}
  \begin{columns}[T]
    \begin{column}{0.42\textwidth}
      \textbf{Bài toán:} tấn công CIA suy ra topology chỉ qua truy vấn hộp đen;
      lộ topology $\to$ mất tài sản trí tuệ, lộ agent then chốt.
      \vspace{4pt}

      \textbf{Phương pháp: 3 bước}
      \begin{enumerate}
        \item \emph{Phantom topology:} đi dây lại cạnh, độ lệch $\rho \ge \rho_0$ ($k$-anonymity).
        \item \emph{Semantic edge realization:} LLM sinh bằng chứng giả cho phantom edge; paraphrase để giấu cạnh thật.
        \item \emph{Dual-view execution:} chạy trên $G$ thật, chỉ phơi bày $G'$.
      \end{enumerate}
    \end{column}
    \begin{column}{0.56\textwidth}
      \centering
      {\small AUC / ACC / F1 của tấn công (G-Designer, thấp hơn là tốt)\par}
      \vspace{3pt}
      \footnotesize
      \begin{tabular}{@{}l c c@{}}
        \toprule
        Dataset & Delimiters & MIRAGE \\
        \midrule
        MMLU      & 0.76 / 0.61 / 0.64 & \textbf{0.67 / 0.50 / 0.49} \\
        GSM8K     & 0.80 / 0.76 / 0.67 & \textbf{0.64 / 0.51 / 0.42} \\
        SVAMP     & 0.70 / 0.68 / 0.71 & \textbf{0.64 / 0.48 / 0.39} \\
        HumanEval & 0.76 / 0.65 / 0.56 & \textbf{0.68 / 0.44 / 0.52} \\
        \bottomrule
      \end{tabular}
      \vspace{6pt}
      \begin{alertblock}{Hạn chế}
        \raggedright
        Chỉ 1 tấn công, không có adaptive attacker; AUC vẫn 0.64--0.68 (lý tưởng 0.5);
        không báo cáo chi phí gọi LLM thêm.
      \end{alertblock}
    \end{column}
  \end{columns}
\end{frame}

\begin{frame}{TIDE: misinformation làm chệch collective fact recovery}
  \textbf{Thiết lập}
  \begin{itemize}
    \item 120 môi trường Hi-ToM; 5 agent, mỗi agent chỉ thấy một phần chuỗi vị trí.
    \item C0: tất cả trung thực · C1: key agent nói dối (Event / Coverage / Source).
    \item GPT-5.5, DeepSeek-v4-pro, Grok-4.5; đo cả evidence lineage.
  \end{itemize}
  \vspace{4pt}
  \textbf{Kết quả chính}
  \begin{itemize}
    \item Tỉ lệ đa số đúng: \textbf{72.50\% $\to$ 14.17\%} (C0 $\to$ C1).
    \item Testimony sai được chấp nhận \textbf{70.7\%}, so với \textbf{25.4\%} của testimony đúng.
    \item Propagation order: \textbf{2.66} (sai) so với \textbf{1.33} (đúng).
  \end{itemize}
  \begin{block}{Thông điệp}
    Thảo luận \emph{khuếch đại} chiều hướng của bằng chứng ban đầu thay vì sửa sai;
    $\sim$85\% khoảng cách hình thành ngay khi chia sẻ testimony.
    Observer giảm 90\% decoy majority nhưng không tăng đáp án đúng.
    Hạn chế: 1 tác vụ hẹp, chưa có defense.
  \end{block}
\end{frame}

% ============================================================
\section{Orchestration \& Control}
% ============================================================

\begin{frame}{ProgRouter: Đặt vấn đề}
  \textbf{Bối cảnh}
  \begin{itemize}
    \item Multi-agent LLM workflow giải tốt task phức tạp nhưng \alert{đắt}:
      gọi LLM nhiều lần, context tích lũy dài $\to$ tốn token, năng lượng GPU, latency.
    \item Cascade routing (RouteLLM, CASCADIA, MasRouter) chọn model
      \emph{một lần cho cả query} -- không theo kịp trạng thái thay đổi qua từng bước.
  \end{itemize}
  \vspace{6pt}
  \begin{block}{Câu hỏi của paper}
    Chọn LLM cho \textbf{từng bước} của workflow thế nào để giữ chất lượng
    mà vẫn nằm trong ngân sách chi phí?
  \end{block}
\end{frame}

\begin{frame}{ProgRouter: Problem setting}
  \small
  \begin{columns}[T]
    \begin{column}{0.46\textwidth}
      \textbf{Workflow} $w = \{C, R, M\}$
      \begin{itemize}
        \item $C$: coordinator agent -- chia subtask, giữ \emph{ledger} trạng thái $s_t$,
          chọn role $r_t$ cho bước tiếp theo.
        \item $R$: các role worker (analyst, coder…).
        \item $M$: model zoo nhiều cỡ; $M_r \subseteq M$ là các model dùng được cho role $r$.
      \end{itemize}
    \end{column}
    \begin{column}{0.5\textwidth}
      \textbf{Bài toán tối ưu online} trên dòng $W$ task:
      tìm policy $\pi(s_t, r_t) \mapsto m_t \in M_{r_t}$
      \begin{align*}
        \max_\pi\;& \tfrac{1}{W}\textstyle\sum_{w} P(w) \\
        \text{s.t.}\;& \textstyle\sum_{t} E(m_t) \le E_w,\quad
                       \sum_{t} \Gamma(m_t) \le \Gamma_w \quad \forall w \\
        & \tfrac{1}{W}\textstyle\sum_{w}\sum_{t} E(m_t) \le \tilde E
      \end{align*}
      {\footnotesize $P(w)\in[0,1]$: chất lượng cuối; $E$, $\Gamma$: năng lượng, thời gian;
      $\tilde E$: ngân sách năng lượng trung bình dài hạn.}
      \begin{alertblock}{Hai khó khăn}
        (1) $P(w)$ chỉ biết khi workflow kết thúc $\to$ không chấm được từng bước.\\
        (2) Ràng buộc trung bình phụ thuộc các task tương lai chưa biết.
      \end{alertblock}
    \end{column}
  \end{columns}
\end{frame}

\begin{frame}{ProgRouter: Tổng quan phương pháp}
  \centering
  \begin{tikzpicture}[node distance=0.42cm,
      box/.style={stage, text width=2.1cm, font=\small, minimum height=1.7cm},
      core/.style={box, fill=mainblue, text=white}]
    \node[box] (c) {\textbf{Coordinator}\\ ledger $s_t$,\\ chọn role $r_t$};
    \node[core, right=of c] (g) {\textbf{Scorer}\\ tiến độ hiện tại\\ $g(s_t)$};
    \node[core, right=of g] (p) {\textbf{Predictor}\\ mức tăng $\hat y_t(m)$\\ $\forall m \in M_{r_t}$};
    \node[core, right=of p] (r) {\textbf{Router}\\ $m^\star_t = \arg\max$\\ score$(m)$};
    \node[box, right=of r] (x) {\textbf{Worker} $r_t$\\ chạy bằng $m^\star_t$\\ $\to s_{t+1}$};
    \foreach \a/\b in {c/g,g/p,p/r,r/x} \draw[arrow] (\a) -- (\b);
    \draw[arrow] (x.south) -- ++(0,-1.05) -| node[pos=0.25, below, font=\footnotesize]
      {trạng thái mới $s_{t+1}$ quay lại coordinator} (c.south);
    \draw[arrow, dashed] (x.north) -- ++(0,0.6) -| node[pos=0.25, above, font=\footnotesize]
      {nhãn học online: $y_t = g(s_{t+1}) - g(s_t)$} (p.north);
    \node[draw=mainblue, dotted, thick, rounded corners, fit=(g)(p)(r), inner sep=4pt,
      label={[font=\footnotesize\bfseries, text=mainblue]below:ProgRouter}] {};
  \end{tikzpicture}
  \vspace{10pt}

  Ở mỗi bước, coordinator chọn role; ProgRouter \textbf{đo} tiến độ hiện tại,
  \textbf{dự đoán} mỗi model ứng viên giúp tiến thêm bao nhiêu, rồi \textbf{chọn} model
  cân bằng giữa tiến độ và ngân sách.
\end{frame}

\begin{frame}{ProgRouter: Đo và dự đoán tiến độ}
  \begin{columns}[T]
    \begin{column}{0.53\textwidth}
      \begin{block}{Multi-View Task Progress Representation}
        \small
        \[ g(s_t) = b + \alpha\, r + \beta\, d + \gamma\, e \;\in [0,1] \]
        \vspace{-10pt}
        \begin{itemize}\setlength{\itemsep}{1pt}
          \item $b$: \textbf{Overall outcome} -- Trạng thái tổng quát của workflow.
          \item $r$: \textbf{Subtask completion} -- Tỉ lệ subtask hoàn thành.
          \item $d$: \textbf{Progress trend} -- Tiến độ của workflow.
          \item $e$: \textbf{State quality} -- Mức ý nghĩa về ngữ nghĩa và cấu trúc
            của lần chuyển trạng thái gần nhất.
        \end{itemize}
      \end{block}
    \end{column}
    \begin{column}{0.45\textwidth}
      \begin{block}{Online Task Progress Predictor}
        \centering
        \begin{tikzpicture}[
            box/.style={stage, text width=2.35cm, minimum height=0.95cm, font=\scriptsize}]
          \node[box, text width=0.85cm] (s) at (0,0) {$s_t$,\\ $m_t$};
          \node[box] (a) at (2.05, 0.75) {\textbf{Structured path}\\[2pt] $\hat y^{\text{str}}_t = P_{\text{str}}(x^{\text{str}}_t)$};
          \node[box] (b) at (2.05,-0.75) {\textbf{Semantic path}\\[2pt] $\hat y^{\text{sem}}_t = P_{\text{sem}}(x^{\text{sem}}_t)$};
          \node[box, text width=1.75cm, fill=mainblue, text=white] (c) at (4.75,0) {\textbf{Meta-gated}\\ \textbf{learner}};
          \draw[arrow] (s.east) -- (a.west); \draw[arrow] (s.east) -- (b.west);
          \draw[arrow] (a.east) -- (c.west); \draw[arrow] (b.east) -- (c.west);
          \node[below=0.25cm of c] (y) {$\hat y_t(m_t)$};
          \draw[arrow] (c) -- (y);
        \end{tikzpicture}
        \[ \hat y_t = P_{\text{meta}}\bigl(\hat y^{\text{str}}_t,\; \hat y^{\text{sem}}_t\bigr) \]
      \end{block}
    \end{column}
  \end{columns}
\end{frame}

\begin{frame}{ProgRouter: Chọn model ở mỗi bước}
  Ở bước $t$, chọn $m^\star_t = \arg\max_{m_t \in M_{r_t}} \mathrm{score}(m_t)$, với mỗi model ứng viên $m_t$:
  \[ \mathrm{score}(m_t) =
     \underbrace{V(1-g(s_t))\,\hat y_t(m_t)}_{\text{\textcolor{mainblue}{lợi ích}}}
     \;-\; \underbrace{Q_w\bigl(E(m_t)-\tilde E\bigr)}_{\text{\textcolor{mainblue}{phạt dài hạn}}}
     \;-\; \underbrace{c^\Gamma_t\,\Gamma(m_t) \;-\; c^E_t\,E(m_t)}_{\text{\textcolor{mainblue}{phạt trong task}}} \]
  \textbf{\color{mainblue}Lợi ích:} task còn xa đích ($g$ nhỏ) thì tiến độ đáng giá hơn $\to$ dám gọi model mạnh.
  \vspace{2pt}
  \begin{columns}[T]
    \begin{column}{0.48\textwidth}
      \begin{block}{Phạt dài hạn}
        Hàng đợi ảo cộng dồn phần vượt ngân sách sau mỗi task:
        \[ Q_{w+1} = \max\bigl\{0,\; Q_w + \textstyle\sum_{t=1}^{T_w} E(m_t) - \tilde E\bigr\} \]
      \end{block}
    \end{column}
    \begin{column}{0.48\textwidth}
      \begin{block}{Phạt trong task}
        Tăng dần khi task gần hết ngân sách ($c^\Gamma_t$ tương tự, theo thời gian):
        \[ c^E_t = \exp\Bigl(\tfrac{E(m_t) + \sum_{i<t} E(m_i)}{E_w}\Bigr) \]
      \end{block}
    \end{column}
  \end{columns}
\end{frame}

\begin{frame}{ProgRouter: Thực nghiệm và kết quả chính}
  \relax{\footnotesize
  \textbf{Thiết lập:} 4 benchmark (code, toán, RAG QA); model zoo 7--9 LLM từ 0.5B đến 35B;
  coordinator Qwen3-Coder-Next 80B; năng lượng GPU đo thật bằng NVML.\par}
  \vspace{3pt}
  \centering
  \scriptsize
  \begin{tabular}{@{}l c c c c@{}}
    \toprule
    Pass/Precision \% / năng lượng J & HumanEval+ & MBPP & MATH-500 & ASQA \\
    \quad (ngân sách $\tilde E$) & (4,800) & (4,500) & (7,000) & (19,000) \\
    \midrule
    Model lớn nhất, dùng một mình & 97.1 / \alert{5,443} & 93.9 / \alert{4,804} & 92.2 / \alert{26,276} & 92.3 / \alert{21,887} \\
    Educated Guessing & 91.5 / \alert{4,916} & 65.2 / 3,831 & 79.3 / \alert{7,023} & 90.8 / \alert{23,423} \\
    CASCADIA          & 84.8 / 4,658 & 78.5 / 3,857 & \textbf{87.8} / 6,875 & 89.8 / \alert{27,857} \\
    MasRouter         & 90.9 / 4,483 & 67.8 / 3,700 & 73.6 / \alert{8,294} & 89.8 / 16,368 \\
    \textbf{ProgRouter} & \textbf{93.0} / 4,796 & \textbf{79.4} / \textbf{3,376} & 84.3 / \textbf{6,112} & \textbf{92.1} / 18,373 \\
    \bottomrule
  \end{tabular}\\[2pt]
  {\scriptsize \alert{Đỏ}: vượt ngân sách năng lượng dài hạn; \textbf{đậm}: tốt nhất trong các phương pháp không vượt ngân sách.\par}
  \vspace{2pt}
  \raggedright\footnotesize
  \begin{itemize}
    \item Dùng một model cố định đều thua: model lớn vượt ngân sách; model nhỏ cũng tốn vì phải sửa lỗi nhiều lần.
    \item ProgRouter luôn trong ngân sách, tốt nhất ở 3/4 benchmark; phần lớn lời gọi dồn vào model nhỏ
      (84.3\% Qwen2.5-Coder 0.5B trên HumanEval+, 91.0\% Granite 3B trên MATH-500).
  \end{itemize}
\end{frame}


\begin{frame}{HalluProp: Đặt vấn đề}
  \textbf{Bối cảnh}
  \begin{itemize}
    \item Trong MAS, hallucination \alert{lan truyền}: một agent sai, các agent khác
      chấp nhận, lặp lại và khuếch đại $\to$ cả hệ thống ra kết quả sai.
    \item Các phương pháp hiện có (SelfCheckGPT, LLM Debate, Inspector, GUARDIAN) đều
      \textbf{post-hoc}: chạy hết MAS rồi mới phát hiện -- lỗi đã lan, tốn 53--308 s/query.
  \end{itemize}
  \vspace{6pt}
  \begin{block}{Câu hỏi của paper}
    Có thể dự đoán agent nào sẽ hỏng \textbf{trước khi} các agent bắt đầu giao tiếp?
  \end{block}
\end{frame}

\begin{frame}{HalluProp: Problem setting}
  \small
  \begin{columns}[T]
    \begin{column}{0.47\textwidth}
      \textbf{MAS là đồ thị có hướng} $M = (V, E)$
      \begin{itemize}
        \item $V = \{v_1, \dots, v_N\}$: $N$ agent; agent $v_i$ có system prompt (role) $p_i$.
        \item $E \in \{0,1\}^{N \times N}$: cạnh $v_j \to v_i$ nghĩa là $v_i$ đọc output của $v_j$.
        \item $q$: câu hỏi của người dùng.
        \item $H_i \in \{0,1\}$: agent $v_i$ có hallucinate không
          (thực nghiệm: trả lời sai so với đáp án).
      \end{itemize}
    \end{column}
    \begin{column}{0.49\textwidth}
      \textbf{Mục tiêu:} ước lượng rủi ro của từng agent
      \[ h_i = P\bigl(H_i = 1 \mid \{p_i\}_{i=1}^N,\; q,\; E\bigr) \]
      Đầu vào \textbf{chỉ có} role, câu hỏi và topology -- chưa có câu trả lời nào.
      \begin{block}{Tách rủi ro làm hai phần}
        \begin{itemize}
          \item \textbf{Intrinsic} $\pi_i \in [0,1]$: rủi ro agent $v_i$ tự sai.
          \item \textbf{Propagation} $t_{j\to i} \in [0,1]$: xác suất lỗi của $v_j$ lây sang $v_i$.
        \end{itemize}
        Gộp $\pi$ và $T = [t_{j\to i}]$ trên đồ thị để ra $h = [h_1, \dots, h_N]$.
      \end{block}
    \end{column}
  \end{columns}
\end{frame}

\begin{frame}{HalluProp: Intrinsic risk -- role có hợp với câu hỏi không?}
  \small
  \begin{columns}[T]
    \begin{column}{0.55\textwidth}
      \textbf{1. Encode} role và câu hỏi vào cùng không gian:
      \[ e^R_i = \mathrm{Enc}(p_i), \qquad e^Q = \mathrm{Enc}(q) \]
      \textbf{2. Gate theo câu hỏi} -- chọn các chiều của role liên quan đến câu hỏi:
      \[ g_i = \sigma\bigl(W_{\text{gate}}\,(e^R_i \odot e^Q) + b_g\bigr) \]
      \textbf{3. Tách role} làm hai phần:
      \[ \underbrace{\kappa_i = e^R_i \odot g_i}_{\text{competence}}, \qquad
         \underbrace{\rho_i = \mathrm{LN}(e^R_i - \kappa_i)}_{\text{susceptibility}} \]
      \textbf{4. Rủi ro:}
      \[ \pi_i = \sigma\bigl(w_\rho^\top \rho_i - w_\kappa^\top \kappa_i + b_\pi\bigr) \]
    \end{column}
    \begin{column}{0.41\textwidth}
      \begin{block}{Trực giác}
        \begin{itemize}
          \item $\kappa_i$ (\emph{competence}): phần của role khớp với câu hỏi.
          \item $\rho_i$ (\emph{susceptibility}): phần còn lại -- lạc đề, mơ hồ hoặc gây nhiễu.
          \item Nhiều nhiễu, ít năng lực $\to$ $\pi_i$ cao.
        \end{itemize}
      \end{block}
      {\footnotesize Chỉ so độ tương đồng tổng thể giữa role và câu hỏi là không đủ,
      vì role prompt thường mơ hồ hoặc chứa nội dung gây nhiễu.}
    \end{column}
  \end{columns}
\end{frame}

\begin{frame}{HalluProp: Lan truyền và Noisy-OR}
  \small
  \begin{columns}[T]
    \begin{column}{0.48\textwidth}
      \begin{block}{Propagation qua cạnh $v_j \to v_i$}
        \[ t_{j\to i} = \sigma\bigl(\gamma_{j,i}\; \kappa_j^\top W_{\text{sem}}\, \rho_i + b_t\bigr) \]
        \vspace{-8pt}
        \begin{itemize}
          \item $\kappa_j$: bên gửi có sức ảnh hưởng.
          \item $\rho_i$: bên nhận dễ bị lây.
          \item $\gamma_{j,i} = \phi(\hat k_{\text{out}}(j), \hat k_{\text{in}}(i)) \ge 1$:
            khuếch đại theo bậc ra / bậc vào.
          \item $t_{j\to i} = 0$ nếu không có cạnh.
        \end{itemize}
      \end{block}
    \end{column}
    \begin{column}{0.48\textwidth}
      \begin{block}{Noisy-OR trên đồ thị}
        \[ h_i = 1 - (1-\pi_i) \prod_{j \in N_{\text{in}}(i)} \bigl(1 - t_{j\to i}\, h_j\bigr) \]
        Agent \emph{không} hỏng $\iff$ tự không sai \textbf{và} không bị lây từ agent nào.
      \end{block}
    \end{column}
  \end{columns}
  \vspace{4pt}
  \textbf{Huấn luyện 2 giai đoạn:} học $\pi$ với nhãn từng agent \emph{trước} giao tiếp,
  rồi học $T$ với nhãn \emph{sau} giao tiếp, BCE trên $h$.
\end{frame}

\begin{frame}{HalluProp: Thực nghiệm và kết quả}
  \begin{columns}[T]
    \begin{column}{0.46\textwidth}
      \textbf{Thiết lập}
      \begin{itemize}\small
        \item MMLU, MATH, MedQA (500 câu mỗi bộ).
        \item 4 backbone: GPT-3.5, GPT-4o, Llama-3-8B, Qwen-2-7B.
        \item Topology ngẫu nhiên; role sạch lẫn role bị nhiễu; 3 vòng giao tiếp.
      \end{itemize}
      \textbf{Phát hiện agent hỏng}
      \begin{itemize}\small
        \item AUROC trung bình \textbf{84.6}; Hit@1 \textbf{83.9}; Recall@3 \textbf{76.2}.
      \end{itemize}
      \textbf{Tốc độ}
      \begin{itemize}\small
        \item Dưới \textbf{1 s}, so với 53--308 s của post-hoc (nhanh hơn $>65\times$).
      \end{itemize}
    \end{column}
    \begin{column}{0.5\textwidth}
      \centering
      {\small Accuracy sau can thiệp vào agent rủi ro nhất\\ (trung bình 4 LLM $\times$ 3 dataset)\par}
      \vspace{4pt}
      \small
      \begin{tabular}{@{}l c@{}}
        \toprule
        Phương pháp & Accuracy \% \\
        \midrule
        LLM Debate   & 55.7 \\
        Inspector    & 54.4 \\
        SelfCheckGPT & 56.7 \\
        GUARDIAN     & 63.5 \\
        \midrule
        HalluProp-Pruning     & 63.6 \\
        HalluProp-Refinement  & \textbf{67.6} \\
        \bottomrule
      \end{tabular}\\[4pt]
      \raggedright\footnotesize
      \emph{Pruning}: cắt mọi cạnh ra của agent đó.
      \emph{Refinement}: đổi role của nó thành verifier thận trọng.
      Các baseline đều post-hoc, phải chạy hết MAS.
    \end{column}
  \end{columns}
\end{frame}

\begin{frame}{HalluProp: Ablation}
  \begin{columns}[T]
    \begin{column}{0.5\textwidth}
      \centering
      {\small AUROC trung bình trên GPT-3.5-turbo (3 dataset)\par}
      \vspace{4pt}
      \begin{tabular}{@{}l c c@{}}
        \toprule
        Biến thể & AUROC & Giảm \\
        \midrule
        HalluProp đầy đủ        & \textbf{84.8} & -- \\
        bỏ intrinsic risk       & 61.2 & \alert{$-23.6$} \\
        bỏ propagation          & 66.9 & \alert{$-17.9$} \\
        bỏ topology ($\gamma$)  & 72.1 & \alert{$-12.7$} \\
        \bottomrule
      \end{tabular}
    \end{column}
    \begin{column}{0.46\textwidth}
      \begin{itemize}
        \item \textbf{Bỏ intrinsic risk} (thay bằng độ tương đồng role--câu hỏi): giảm mạnh nhất
          $\to$ tách competence / susceptibility là cần thiết.
        \item \textbf{Bỏ propagation} ($t_{j\to i} = 0$, khi đó $h_i = \pi_i$): lỗi trong MAS
          không chỉ là lỗi riêng lẻ mà còn lan qua giao tiếp.
        \item \textbf{Bỏ topology}: không nhận ra các agent phát tán mạnh.
      \end{itemize}
    \end{column}
  \end{columns}
  \vspace{6pt}
  \begin{block}{}
    Cả ba thành phần đều cần thiết; thiếu bất kỳ thành phần nào, AUROC giảm từ 12.7 đến 23.6 điểm.
  \end{block}
\end{frame}

% ============================================================
\section{Communication \& Interaction}
% ============================================================

\begin{frame}{DAIN: Dynamic Agent-based Interaction Network}
  \centering
  \begin{tikzpicture}[node distance=0.5cm]
    \node[stage, text width=2.8cm] (a) {\textbf{Encode}\\ ResNet-50, BERT $\to e_c$};
    \node[stage, text width=2.8cm, right=of a] (b) {\textbf{Meta-Controller}\\ sparse gate + soft graph $\Phi$};
    \node[stage, text width=2.8cm, right=of b] (c) {\textbf{Message passing}\\ low-rank bottleneck};
    \node[stage, text width=2.8cm, right=of c] (d) {\textbf{Consensus}\\ attention $\to$ kết quả};
    \foreach \x/\y in {a/b,b/c,c/d} \draw[arrow] (\x) -- (\y);
  \end{tikzpicture}
  \raggedright
  \begin{columns}[T]
    \begin{column}{0.48\textwidth}
      \begin{itemize}
        \item 8 agent theo PID: 3 synergy · 3 uniqueness · 2 redundancy.
        \item Loss phạt kích hoạt và mật độ giao tiếp: $\lambda\lVert g\rVert_1 + \lambda\lVert \Phi\rVert^2$.
        \item Chỉ \textbf{3.2--5.5 / 8} agent được kích hoạt.
        \item ADNI accuracy \textbf{73.5} so với 70.9 (best baseline).
      \end{itemize}
    \end{column}
    \begin{column}{0.48\textwidth}
      \begin{alertblock}{Lưu ý khi đánh giá}
        \begin{itemize}
          \item “Agent” là neural module, không phải LLM: bản chất là MoE có message passing.
          \item Bỏ communication: $-1.4$; bỏ dynamic scheduling: $-1.7$ điểm (ADNI).
          \item Số liệu không nhất quán giữa các bảng; baseline yếu.
        \end{itemize}
      \end{alertblock}
    \end{column}
  \end{columns}
\end{frame}

\begin{frame}{AGENT-RADAR: attention steering theo context relevance}
  \textbf{Ý tưởng}
  \begin{itemize}
    \item Giữ toàn bộ lịch sử, không tóm tắt hay cắt bớt.
    \item Chấm điểm từng câu; câu có điểm $\ge \theta = 0.65$ được nhấn mạnh bằng SPA.
    \item Training-free, cắm được vào GPTSwarm, AutoGen, debate.
  \end{itemize}
  \begin{block}{Điểm của câu $s_k$}
    \[ \underbrace{\lambda_s^{\,d_G - 1}}_{\text{spatial decay}} \times
       \underbrace{\lambda_t^{\,t-\tau-1}}_{\text{temporal decay}} \times
       \underbrace{\cos(s_k, q)}_{\text{semantic relevance}} \]
  \end{block}
  \textbf{Hạn chế:} latency +23\% · cần white-box (logit) $\to$ không dùng với API model ·
  tune trên MuSiQue (cũng là test set).
\end{frame}

% ============================================================
\section{Domain Applications}
% ============================================================

\begin{frame}{MAC-RRG: multi-agent cho radiology report generation}
  \centering
  \begin{tikzpicture}[node distance=0.3cm, stage/.append style={text width=2.1cm}]
    \node[stage] (a) {\textbf{X-ray}\\ ảnh ngực};
    \node[stage, right=of a] (b) {\textbf{Draft}\\ Swin + Llama2-7B};
    \node[stage, right=of b] (c) {\textbf{Entities}\\ giải phẫu, bệnh lý};
    \node[stage, right=of c] (d) {\textbf{2 agent}\\ MM-KG + Knowledge};
    \node[stage, right=of d] (e) {\textbf{Report}\\ refined};
    \foreach \x/\y in {a/b,b/c,c/d,d/e} \draw[arrow] (\x) -- (\y);
  \end{tikzpicture}
  \raggedright
  \begin{columns}[T]
    \begin{column}{0.44\textwidth}
      \textbf{Kết quả}
      \begin{itemize}
        \item BLEU-4: \textbf{0.197 / 0.128 / 0.105} trên IU X-Ray / MIMIC-CXR / CheXpert Plus.
        \item Clinical F1 (CheXpert Plus): \textbf{0.295} so với 0.260 của R2GenGPT.
      \end{itemize}
    \end{column}
    \begin{column}{0.52\textwidth}
      \begin{alertblock}{Đánh giá phản biện}
        \begin{itemize}
          \item “Agent” = module retrieval tất định, không giao tiếp hay thương lượng.
          \item “Iterative” thực tế chỉ 1 lượt tinh chỉnh.
          \item BLEU tăng nhỏ (0.101 $\to$ 0.105); không kiểm định thống kê, không có bác sĩ đánh giá.
        \end{itemize}
      \end{alertblock}
    \end{column}
  \end{columns}
\end{frame}

\begin{frame}{UrbanAgent: collaborative reasoning cho urban region profiling}
  \begin{columns}[T]
    \begin{column}{0.48\textwidth}
      \textbf{Stage A · Tool-augmented agents}
      \begin{itemize}
        \item 4 agent theo modality: satellite, text, POI, 3D.
        \item Qwen3-VL-4B + 4 tool (web search, crop-zoom, nightlight, ảnh lịch sử).
        \item SFT $\to$ GRPO, reward = outcome + format + process.
      \end{itemize}
    \end{column}
    \begin{column}{0.48\textwidth}
      \textbf{Stage B · Relation-aware graph}
      \begin{itemize}
        \item \emph{Corroboration:} dự đoán gần nhau $\to$ củng cố.
        \item \emph{Rectification:} lệch nhiều $\to$ sửa.
        \item \emph{Supervisory:} agent tự tin hơn hướng dẫn; 5 vòng GRU $\to$ final judge.
      \end{itemize}
    \end{column}
  \end{columns}
  \vspace{4pt}
  \textbf{Kết quả:} $R^2$ \textbf{0.63--0.65} trong miền (carbon, dân số, GDP);
  $\approx$ \textbf{0.60} trên thành phố chưa thấy; \textbf{+8.1\%} so với baseline mạnh nhất.
  \begin{alertblock}{Đánh giá phản biện}
    “Debate” là GNN/GRU được huấn luyện, không phải tranh luận bằng ngôn ngữ.
    Dữ liệu nhỏ (2,000 vị trí); text sinh từ chính ảnh vệ tinh $\to$ các modality không độc lập.
  \end{alertblock}
\end{frame}

% ============================================================
\section{So sánh \& kết luận}
% ============================================================

\begin{frame}{So sánh hệ thống tám paper}
  \centering
  \footnotesize
  \begin{tabular}{@{}l l c c l c@{}}
    \toprule
    Paper & Loại đóng góp & Training & White-box & Chi phí & Thống kê \\
    \midrule
    MIRAGE      & Defense          & \no    & \no  & \no              & \no \\
    TIDE        & Benchmark        & \no    & \no  & \no              & \yes \\
    ProgRouter  & Router           & online & \no  & \yes{} năng lượng & \no \\
    HalluProp   & Risk predictor   & \yes   & \no  & $\sim$ thời gian & \no \\
    DAIN        & Kiến trúc fusion & \yes   & --   & $\sim$ tham số   & \no \\
    AGENT-RADAR & Plug-in          & \no    & \yes & $\sim$ thời gian & \no \\
    MAC-RRG     & Ứng dụng         & \yes   & --   & \no              & \no \\
    UrbanAgent  & Ứng dụng         & \yes   & --   & \no              & \no \\
    \bottomrule
  \end{tabular}
  \vspace{6pt}

  \raggedright\small
  \textbf{Đọc bảng:} chỉ TIDE có kiểm định thống kê; chỉ ProgRouter đo năng lượng GPU thực;
  chỉ AGENT-RADAR bắt buộc white-box. Cùng benchmark nhưng khác tập con, topology, backbone
  $\to$ không so trực tiếp được con số.
\end{frame}

\begin{frame}{Ba nhận định xuyên suốt}
  \begin{enumerate}
    \item \textbf{Communication graph là trung tâm.}\\
      6/8 bài thao tác trực tiếp: giấu (MIRAGE), phân tích (HalluProp), đo khoảng cách (RADAR),
      học (DAIN, UrbanAgent), truy vết (TIDE).
    \item \textbf{Lan truyền là rủi ro cốt lõi.}\\
      Thảo luận khuếch đại sai lệch thay vì sửa sai (TIDE); hallucination lan theo cạnh (HalluProp).
    \item \textbf{Khái niệm “agent” bị nới rộng.}\\
      LLM agent thật (5 bài) · LLM + aggregator học được (UrbanAgent) ·
      neural / retrieval module (DAIN, MAC-RRG).
  \end{enumerate}
  \vspace{6pt}
  \begin{block}{}
    Hiệu quả không đến từ nhiều agent hơn: ProgRouter dùng chủ yếu model 0.5B;
    DAIN tốt nhất với $\tau = 4/8$ agent.
  \end{block}
\end{frame}

\begin{frame}{Research gap và hướng phát triển}
  \begin{columns}[T]
    \begin{column}{0.48\textwidth}
      \textbf{Research gap}
      \begin{itemize}
        \item[G1] Chưa có benchmark chung: chất lượng + chi phí + robustness.
        \item[G2] 7/8 bài không có phương sai hay kiểm định thống kê.
        \item[G3] Quy mô nhỏ: $\le 8$ agent, 2--5 vòng, tác vụ ngắn.
        \item[G4] Chưa có adaptive attacker; chưa có cách khôi phục sự thật.
        \item[G5] Phương pháp tối ưu chưa được thử trong điều kiện đối kháng.
      \end{itemize}
    \end{column}
    \begin{column}{0.48\textwidth}
      \textbf{Hướng phát triển}
      \begin{itemize}
        \item \emph{Benchmark 3 trục:} chất lượng · chi phí · robustness (giao thức C0/C1 của TIDE).
        \item \emph{Provenance-aware communication:}
          $\mathrm{score}(\text{RADAR}) \times (1 - \hat h_j) \times \mu^{\text{hop}}$.
        \item \emph{Risk-aware routing:} ProgRouter + HalluProp, bước rủi ro $\to$ model mạnh.
        \item \emph{Verifier agent chủ động:} đòi bằng chứng trực tiếp trước khi đồng thuận.
      \end{itemize}
    \end{column}
  \end{columns}
\end{frame}

\begin{frame}{Kết luận}
  \begin{enumerate}
    \item \textbf{Communication graph là trung tâm} -- quyết định hiệu năng, lan truyền rủi ro, cần được bảo vệ.
    \item \textbf{Lan truyền là rủi ro cốt lõi} -- thảo luận khuếch đại sai lệch thay vì sửa sai.
    \item \textbf{Hiệu quả = kích hoạt đúng} -- đúng agent, đúng model, chú ý đúng nội dung.
    \item \textbf{Thực nghiệm còn yếu} -- quy mô nhỏ, thiếu kiểm định, thiếu đo chi phí.
  \end{enumerate}
  \vspace{6pt}
  \begin{block}{Câu hỏi của lĩnh vực đang chuyển}
    Từ \emph{“nhiều agent có tốt hơn một agent?”}\\
    sang \textbf{“làm sao để MAS hiệu quả, đáng tin cậy và an toàn?”}
  \end{block}
\end{frame}

\begin{frame}{Phân công}
  \centering
  \begin{tabular}{@{}l l@{}}
    \toprule
    Thành viên & Paper phụ trách \\
    \midrule
    Trần Xuân Hương & MIRAGE · TIDE \\
    Vũ Tùng Dương   & ProgRouter · HalluProp \\
    Phạm Việt Dũng  & DAIN · AGENT-RADAR \\
    Lê Văn Đức      & MAC-RRG · UrbanAgent \\
    \bottomrule
  \end{tabular}
\end{frame}

{\HUSTUseBackground{theme_hust_oneside}
\begin{frame}[plain]
  \placecontent{0.355\paperwidth}{0.300\paperheight}{0.640\paperwidth}{
    \centering
    {\color{HUSTRed}\bfseries\fontsize{28pt}{36pt}\selectfont CẢM ƠN\\ĐÃ LẮNG NGHE!\par}
    \vspace{14pt}
    {\color{HUSTBlue}\Large Q\&A\par}
  }
\end{frame}
}

\end{document}
